
\documentclass[12pt,a4paper]{article}

\usepackage[a4paper,margin=1in]{geometry}
\usepackage{graphicx}
\usepackage{booktabs}
\usepackage{amsmath,amssymb}
\usepackage{float}
\usepackage{hyperref}
\usepackage{xcolor}
\usepackage{listings}
\usepackage{enumitem}
\usepackage{fancyhdr}
\usepackage{longtable}

\hypersetup{colorlinks=true,linkcolor=blue,urlcolor=blue}

\pagestyle{fancy}
\fancyhf{}
\lhead{ICS1512}
\rhead{Machine Learning Laboratory}
\cfoot{\thepage}

\lstset{
language=Python,
basicstyle=\ttfamily\small,
numbers=left,
frame=single,
breaklines=true,
showstringspaces=false
}

\begin{document}

\begin{tabular}{|p{4cm}|p{11cm}|}
\hline
Experiment No. & 7 \\ \hline
Experiment Title & Bagging, Boosting, and Stacked Ensemble Models
\footnote{\textbf{GitHub Repository: }\url{PASTE_YOUR_GITHUB_LINK_HERE}}\\ \hline
Student Name & Samyuktha V \\ \hline
Register Number & 3122247001053 \\ \hline
Faculty & Dr. Poreddy Ajay Kumar Reddy \\ \hline
Submission Date & 12/08/26 \\ \hline
\end{tabular}

\section{Objective}
To understand ensemble learning strategies -- Bagging, Boosting, and Stacking --
implement a Bagging classifier and Boosting classifiers (AdaBoost and Gradient
Boosting), build a Stacked Ensemble using multiple heterogeneous base learners,
compare all three in terms of accuracy, stability, and generalization, and analyze
the effect of ensemble methods on bias and variance.

\section{Problem Statement}
Given 30 numerical features extracted from digitized images of breast mass
fine-needle aspirates, classify each sample as Malignant or Benign using three
different ensemble strategies, and compare their predictive performance.

\section{Dataset Description}
\begin{longtable}{|p{6cm}|p{8cm}|}
\hline
Dataset Name & Wisconsin Diagnostic Breast Cancer Dataset \\ \hline
Dataset Source & UCI Machine Learning Repository (also available via
\texttt{sklearn.datasets.load\_breast\_cancer()}) \\ \hline
Number of Samples & 569 \\ \hline
Number of Features & 30 numerical attributes \\ \hline
Number of Classes & 2 (Malignant, Benign) \\ \hline
Missing Values & 0 \\ \hline
Train-Test Split & 80\% train (455 samples) / 20\% test (114 samples), stratified,
\texttt{random\_state=42} \\ \hline
\end{longtable}

\section{Exploratory Data Analysis}
\begin{itemize}
\item Dataset overview and missing-value check (none found)
\item Class distribution (357 Benign, 212 Malignant)
\item Correlation heatmap across all 30 features
\item Feature distribution histograms
\end{itemize}

\textbf{Inference:} the dataset is moderately imbalanced (about 63\% Benign / 37\%
Malignant) but clean, with no missing values requiring imputation. Several
features (e.g. \texttt{mean radius}, \texttt{mean perimeter}, \texttt{mean area})
are strongly correlated with each other, as expected since they all describe cell
size from different angles -- this redundancy is exactly the kind of setting where
ensemble methods (which combine many models on different feature/sample subsets)
tend to perform well.

\section{Data Preprocessing}
\begin{itemize}
\item \textbf{Missing values:} none present, no imputation needed.
\item \textbf{Feature scaling:} all 30 features were standardized using
\texttt{StandardScaler}, fit on the training set only, since the Stacked
Ensemble's SVM base learner is distance-based and sensitive to feature scale.
\item \textbf{Train-test split:} 80/20, stratified on the target to preserve the
Malignant/Benign ratio in both sets.
\end{itemize}

\section{Machine Learning Algorithm}

\subsection*{Bagging (Bootstrap Aggregation)}
\textbf{Working principle:} trains multiple Decision Trees, each on a different
bootstrap sample (random sample with replacement) of the training data, and
aggregates their predictions by majority vote.\\
\textbf{Advantages:} reduces variance; trees are trained independently and in
parallel; robust to overfitting compared to a single deep tree.\\
\textbf{Limitations:} does not reduce bias -- if the base learner is systematically
wrong, bagging alone won't fix it.

\subsection*{Boosting (AdaBoost / Gradient Boosting)}
\textbf{Working principle:} trains models sequentially, with each new model
focused on correcting the errors of the previous ones -- AdaBoost reweights
misclassified samples, Gradient Boosting fits each new model to the residual
errors of the ensemble so far.\\
\textbf{Advantages:} reduces bias; often achieves higher accuracy than bagging on
the same base learner.\\
\textbf{Limitations:} sequential training means it can't be parallelized as easily;
more prone to overfitting on noisy data since it keeps focusing on hard examples.

\subsection*{Stacked Ensemble}
\textbf{Working principle:} trains several heterogeneous base models (SVM, Naive
Bayes, Decision Tree), then trains a meta-learner (Logistic Regression) on the
base models' predictions to learn the optimal way to combine them.\\
\textbf{Advantages:} can outperform any individual base model by exploiting their
diverse error patterns; flexible -- any classifier can be a base model or
meta-learner.\\
\textbf{Limitations:} more complex and computationally expensive to train; harder
to interpret than a single model.

\section{Experimental Setup}
\begin{tabular}{ll}
\toprule
Python Version & 3.11 \\
Scikit-Learn Version & 1.9.0 \\
IDE & Jupyter Notebook \\
Operating System & Windows \\
Hardware & Intel i5, 16GB RAM \\
\bottomrule
\end{tabular}

\section{Hyperparameter Selection}

\begin{longtable}{lc}
\toprule
Parameter & Search Space\\
\midrule
Bagging \texttt{n\_estimators} & 10, 50, 100 \\
Bagging \texttt{max\_samples} & 0.5, 0.7, 1.0 \\
Bagging \texttt{max\_features} & 0.5, 0.7, 1.0 \\
Boosting \texttt{n\_estimators} & 50, 100, 150 \\
Boosting \texttt{learning\_rate} & 0.01, 0.1, 1.0 \\
Gradient Boosting \texttt{max\_depth} & 2, 3, 4 \\
Cross-Validation Folds & 5 \\
Search Method & GridSearchCV \\
\bottomrule
\end{longtable}

\section{Results}

\subsection*{Table 1: Bagging Hyperparameter Evaluation (top configurations)}
\begin{longtable}{cccc}
\toprule
n\_estimators & max\_samples & max\_features & Avg CV Accuracy (\%)\\
\midrule
100 & 1.0 & 0.5 & 96.04 \\
100 & 0.7 & 0.5 & 95.82 \\
50  & 1.0 & 0.5 & 95.82 \\
50  & 1.0 & 0.7 & 95.82 \\
100 & 1.0 & 0.7 & 95.82 \\
50  & 0.5 & 0.5 & 95.60 \\
100 & 0.5 & 0.7 & 95.60 \\
50  & 0.7 & 0.7 & 95.38 \\
\bottomrule
\end{longtable}
\textit{Best: n\_estimators=100, max\_samples=1.0, max\_features=0.5 (Avg CV Accuracy
96.04\%, Avg CV F1 0.9684)}

\subsection*{Table 2: Boosting Hyperparameter Evaluation}
\begin{longtable}{lccccc}
\toprule
Model & n\_estimators & learning\_rate & max\_depth & Avg CV Accuracy (\%) & Avg CV F1\\
\midrule
AdaBoost & 100 & 1.0 & -- & 97.36 & 0.9792 \\
\textbf{Gradient Boosting} & 100 & 1.0 & 2 & \textbf{97.58} & \textbf{0.9808} \\
\bottomrule
\end{longtable}
\textit{Gradient Boosting was selected as the best boosting variant and carried
forward to final evaluation.}

\subsection*{Table 3: Stacked Ensemble Evaluation}
\begin{longtable}{lccc}
\toprule
Base Models & Meta Learner & Avg CV Accuracy (\%) & Avg CV F1 Score\\
\midrule
SVM, Naive Bayes, Decision Tree & Logistic Regression & 96.70 & 0.9738 \\
\bottomrule
\end{longtable}

\subsection*{Table 4: Performance Comparison of Ensemble Models (Test Set)}
\begin{longtable}{lcccc}
\toprule
Model & Accuracy (\%) & Precision & Recall & F1 Score\\
\midrule
Bagging & 94.74 & 0.9583 & 0.9583 & 0.9583 \\
Boosting (Gradient Boosting) & 96.49 & 0.9595 & 0.9861 & 0.9726 \\
\textbf{Stacked Ensemble} & \textbf{97.37} & \textbf{0.9726} & \textbf{0.9861} & \textbf{0.9793} \\
\bottomrule
\end{longtable}

\textit{AUC (test set): Bagging = 0.9926, Boosting (Gradient Boosting) = 0.9931,
Stacked Ensemble = 0.9934.}

\section{Result Visualization}

Include, from the notebook: class distribution, correlation heatmap, feature
histograms, confusion matrices for all three final models, the overlaid ROC curve
plot, and the train-vs-test accuracy bar chart.

\textbf{Inference (ROC/AUC):} all three models achieve AUC above 0.99, meaning each
one separates malignant from benign cases very well across nearly all decision
thresholds -- the differences between them are concentrated in the small number of
borderline cases near the default 0.5 threshold, which is exactly where the final
Accuracy/Precision/Recall/F1 numbers in Table 4 differ.

\textbf{Inference (Bias-Variance chart):} Bagging and Boosting both reach 100\%
training accuracy but drop to 94.74\% and 96.49\% respectively on the test set --
a noticeable train-test gap indicating some overfitting, more pronounced for
Bagging. The Stacked Ensemble shows a smaller gap (98.46\% train vs 97.37\% test),
suggesting its meta-learner is generalizing better rather than memorizing the
training set.

\section{Discussion}

\subsection*{How does Bagging reduce variance?}
By training many Decision Trees on different bootstrap samples and averaging
their votes, the individual trees' overfitting mistakes (which differ from tree to
tree since each sees different data) tend to cancel out in the aggregate,
producing a smoother, less variance-prone decision boundary than any single deep
tree.

\subsection*{How does Boosting address model bias?}
Each new model in the boosting sequence is trained specifically to correct the
errors of the previous ensemble (via reweighting in AdaBoost, or fitting to
residuals in Gradient Boosting). This lets a sequence of otherwise weak/biased
learners combine into a much stronger overall model, directly attacking
underfitting rather than just averaging over variance.

\subsection*{Why does stacking benefit from heterogeneous models?}
SVM, Naive Bayes, and Decision Tree make errors on different examples because
they rely on different assumptions (margin maximization, feature independence,
axis-aligned splits, respectively). A meta-learner trained on their combined
outputs can learn which base model to trust more in which situation, capturing
patterns no single base model sees alone -- which is reflected here in the
Stacked Ensemble outperforming both individual ensemble strategies on the test
set.

\subsection*{Which ensemble method performed best and why?}
The \textbf{Stacked Ensemble} performed best overall (97.37\% test accuracy, F1 =
0.9793, AUC = 0.9934), narrowly ahead of Gradient Boosting and more clearly ahead
of Bagging. This matches the bias-variance analysis above: Stacking combines
Bagging-style diversity (heterogeneous base models) with Boosting-style
correction (a meta-learner that weighs each base model's strengths), while
showing the smallest train-test accuracy gap of the three -- indicating the best
balance between fitting the data and generalizing to unseen samples.

\section{Conclusion}
All three ensemble strategies substantially outperform what a single base
classifier would typically achieve on this dataset, confirming the core premise
of ensemble learning. Bagging reduced variance at the cost of some residual
overfitting; Boosting (specifically Gradient Boosting) reduced bias and achieved
higher test accuracy than Bagging; and the Stacked Ensemble, by combining diverse
base learners under a meta-learner, achieved the best overall test performance
and the most balanced bias-variance trade-off of the three.

\section{References}
\begin{enumerate}
\item Scikit-learn: Ensemble Methods, \url{https://scikit-learn.org/stable/modules/ensemble.html}
\item Scikit-learn: Bagging Classifier, \url{https://scikit-learn.org/stable/modules/generated/sklearn.ensemble.BaggingClassifier.html}
\item UCI Machine Learning Repository: Breast Cancer Wisconsin (Diagnostic) Dataset
\end{enumerate}

\end{document}
